% =============================================================================
% CAPÍTULO 4 — DISEÑO DETALLADO DEL BACKEND (EthosBack)
% Archivo maestro del capítulo. Importa las secciones desde sections/cap04/.
% =============================================================================
\chapter{Diseño Detallado del Backend (EthosHub)}
\label{ch:backend}
\resumenCapitulo{Este capítulo desciende al nivel de código del servidor: la arquitectura
de la API REST y sus convenciones, las clases y DTOs de los dominios críticos, las
secuencias de las operaciones clave (autenticación, IA y CV), la comunicación en tiempo
real, la seguridad y el manejo de excepciones, y el empaquetado y despliegue.}

Este capítulo desciende al nivel de \textbf{código} del Modelo C4 para el componente
servidor \texttt{EthosBack}. Mientras el capítulo~\ref{ch:arquitectura} (pág.~\pageref{ch:arquitectura})
describió la arquitectura global, aquí se documenta el diseño interno con la granularidad
que exige la norma ISO/IEC/IEEE~15289:2024 para el ítem de información \textit{Software
Design Description}: la organización de la API REST y sus convenciones; la jerarquía de
clases, servicios y objetos de transferencia de los dominios críticos; las secuencias de
las operaciones de mayor valor (autenticación federada, asistencia de IA, compilación de
currículums); la arquitectura de comunicación en tiempo real; y, por último, la seguridad
del servidor con su modelo centralizado de manejo de excepciones.

Todo el contenido es \textbf{trazable} a paquetes y clases reales del repositorio. El
backend se ejecuta sobre \textbf{Java~17} y \textbf{Spring~Boot~4.0.x}, persiste mediante
\textbf{jOOQ} sobre PostgreSQL~15.10 (servidor TIS) y se empaqueta como un único \textit{JAR}
monolítico que sirve además el \textit{bundle} estático del frontend.

\begin{NotaTecnica}
La organización del código sigue el patrón \textbf{\textit{vertical slice}} (corte
vertical por funcionalidad): cada dominio de negocio agrupa en un mismo paquete su
\texttt{Controller}, su \texttt{Service}, su \texttt{Repository} y sus DTOs. Esto se
contrapone a la organización tradicional por capas técnicas horizontales y reduce el
acoplamiento entre dominios, alineándose con la decisión de arquitectura ADR-03
(véase la sección~\ref{sec:adr}, pág.~\pageref{sec:adr}).
\end{NotaTecnica}

% --- Secciones del capítulo (orden del índice instrucciones.md §4) ---
% =============================================================================
% 4.1 ARQUITECTURA DE LA API REST E INTERFACES
% =============================================================================
\section{Arquitectura de la API REST e Interfaces}
\label{sec:api_rest}

El backend expone una API REST \textbf{stateless} (sin estado de sesión en el servidor):
la identidad del llamante viaja en cada petición como un \textit{JSON Web Token} (JWT)
de tipo \textit{Bearer}, emitido por Supabase y \textbf{verificado} —nunca emitido— por
el servidor. Esta decisión es la piedra angular de la escalabilidad horizontal: cualquier
instancia del \textit{JAR} puede atender cualquier petición sin compartir estado de sesión.

La superficie HTTP se compone de \textbf{25 controladores} agrupados por dominio. Cada
uno delega la lógica en un \texttt{Service} transaccional, que a su vez accede a los datos
mediante un \texttt{Repository} tipado con jOOQ. La figura~\ref{fig:pipeline-http}
(pág.~\pageref{fig:pipeline-http}) anticipa el recorrido completo de una petición a través
de la cadena de filtros y la pila de capas.

\subsection{Pipeline de procesamiento de una petición HTTP}
\label{subsec:pipeline}

Antes de catalogar las rutas conviene comprender el \textbf{ciclo de vida} de una
petición. Toda llamada atraviesa, en orden, la cadena de filtros de servlet, la cadena
de seguridad de Spring Security, el \textit{dispatcher} de Spring MVC y, finalmente, la
pila Controller~$\rightarrow$~Service~$\rightarrow$~Repository. La
figura~\ref{fig:pipeline-http} formaliza este recorrido.

\vspace{0.3cm}
\begin{figure}[h!]
\centering
\begin{tikzpicture}[node distance=0.55cm and 0.9cm,
    etapa/.style={rectangle, rounded corners=3pt, draw=c4Sistema!70!black, fill=c4Componente!30,
        font=\sffamily\scriptsize, align=center, minimum width=2.5cm, minimum height=0.85cm, inner sep=3pt},
    filtro/.style={etapa, fill=colorAcento!25, draw=colorAcento!80!black},
    datos/.style={etapa, fill=c4Datos!22, draw=c4Datos!70!black},
    rel/.style={-{Stealth[length=2mm]}, semithick, draw=grisISO}]
    \node[etapa, fill=c4Persona!25] (cli) {Cliente\\\scriptsize SPA / Axios};
    \node[filtro, right=of cli] (f1) {Filtros servlet\\\scriptsize Cookie · CORS · Headers};
    \node[filtro, right=of f1] (sec) {Spring Security\\\scriptsize JWT Resource Server};
    \node[etapa, below=0.9cm of sec] (ctrl) {\texttt{@RestController}\\\scriptsize valida \& enruta};
    \node[etapa, left=of ctrl] (svc) {\texttt{@Service}\\\scriptsize \texttt{@Transactional}};
    \node[datos, left=of svc] (repo) {\texttt{Repository}\\\scriptsize jOOQ DSLContext};
    \node[datos, below=0.85cm of repo] (db) {PostgreSQL\\\scriptsize esquema \texttt{core}};

    \draw[rel] (cli) -- node[c4etiqueta, above]{\scriptsize HTTP} (f1);
    \draw[rel] (f1) -- (sec);
    \draw[rel] (sec) -- node[c4etiqueta, right]{\scriptsize Authentication} (ctrl);
    \draw[rel] (ctrl) -- node[c4etiqueta, above]{\scriptsize DTO} (svc);
    \draw[rel] (svc) -- node[c4etiqueta, above]{\scriptsize SQL tipado} (repo);
    \draw[rel] (repo) -- (db);
    \draw[rel, dashed] (svc) to[bend left=12] node[c4etiqueta, left]{\scriptsize ApiResponse} (ctrl);
\end{tikzpicture}
\caption{Pipeline de procesamiento de una petición HTTP en el backend.}
\label{fig:pipeline-http}
\end{figure}

\subsection{Enrutamiento global y convenciones de nomenclatura}
\label{subsec:enrutamiento}

Las rutas siguen un esquema \textbf{jerárquico y predecible} que distingue versiones y
ámbitos. La tabla~\ref{tab:rutas-base} enumera las rutas base reales de cada controlador,
extraídas directamente de sus anotaciones \comando{@RequestMapping}. Para evitar
desbordamientos de margen, las rutas largas se presentan con quiebre automático mediante
el formateador de rutas del documento.

\vspace{0.2cm}
\noindent
\renewcommand{\arraystretch}{1.35}
\fontsize{8}{9.6}\selectfont\sffamily
\begin{tabularx}{\textwidth}{|L{4.5cm}|Y|}
\hline
\rowcolor{headerTabla}
\sffamily\textbf{Ruta base} & \sffamily\textbf{Controlador y responsabilidad} \\
\hline
\pathbreak{/api/auth} & \texttt{AuthController}: registro, \textit{login} y sincronización OAuth (\texttt{/oauth/sync}). \\
\hline
\rowcolor{grisTecnico}
\pathbreak{/api/auth/forgot-password} & \texttt{ForgotPasswordController}: solicitud, verificación y reinicio de contraseña vía OTP. \\
\hline
\pathbreak{/api/v1/auth} & \texttt{SecurityController}: cambio de contraseña/correo con OTP, eliminación y exportación de cuenta. \\
\hline
\rowcolor{grisTecnico}
\pathbreak{/api/v1/profile} & \texttt{ProfileController}: perfil profesional básico (bio, foto, datos de contacto). \\
\hline
\pathbreak{/api/projects} · \pathbreak{/api/skills} & \texttt{ProjectController}, \texttt{SkillController}: contenido profesional (rutas a nivel de método). \\
\hline
\rowcolor{grisTecnico}
\pathbreak{/api/v1/work-experiences} & \texttt{WorkExperienceController}: experiencia laboral con geolocalización (lat/lng). \\
\hline
\pathbreak{/api/v1/academic-records} & \texttt{AcademicRecordController}: formación académica. \\
\hline
\rowcolor{grisTecnico}
\pathbreak{/api/v1/portfolio} & \texttt{PortfolioSettingsController}: ajustes del portafolio (SEO, visibilidad, contraseña). \\
\hline
\pathbreak{/api/v1/portfolio/public} & \texttt{PortfolioPublicController}: vista pública del portafolio (sin autenticación). \\
\hline
\rowcolor{grisTecnico}
\pathbreak{/api/v1/cv-documents} & \texttt{CvDocumentController}: CRUD de CV, asistencia IA y compilación a PDF. \\
\hline
\pathbreak{/api/v1/connections} & \texttt{ConnectionController}: conexiones a aplicaciones externas (Drive, etc.). \\
\hline
\rowcolor{grisTecnico}
\pathbreak{/api/v1/messages} & \texttt{ChatController}: envío idempotente de mensajes de chat. \\
\hline
\pathbreak{/api/v1/notifications} & \texttt{NotificationController}: notificaciones del usuario. \\
\hline
\rowcolor{grisTecnico}
\pathbreak{/api/v1/preferences} & \texttt{PreferenceController}: preferencias e idioma del usuario. \\
\hline
\pathbreak{/api/v1/talents} & \texttt{TalentController}: descubrimiento y filtrado de talento. \\
\hline
\rowcolor{grisTecnico}
\pathbreak{/api/v1/recruiter} · \pathbreak{/api/v1/recruiter/profile} & \texttt{RecruiterAiController}, \texttt{RecruiterProfileController}: \textit{insights} IA y perfil de empresa. \\
\hline
\pathbreak{/api/uploads} & \texttt{FileUploadController}: subida de archivos a Supabase Storage. \\
\hline
\rowcolor{grisTecnico}
\pathbreak{/api/admin} · \pathbreak{/api/admin/moderation} · \pathbreak{/api/admin/profiles} · \pathbreak{/api/admin/skills} · \pathbreak{/api/admin/email} & Familia \texttt{Admin*Controller}: métricas, moderación, perfiles, \textit{skills} y correo. \\
\hline
\end{tabularx}
\label{tab:rutas-base}

\begin{NotaTecnica}
Coexisten dos familias de prefijos: las rutas \textbf{legadas} (\pathbreak{/api/auth},
\pathbreak{/api/projects}, \pathbreak{/api/skills}) y las rutas \textbf{versionadas}
(\pathbreak{/api/v1/...}). El versionado explícito en la ruta facilita la evolución
incompatible de contratos sin romper clientes existentes, criterio MAYOR del versionado
semántico (SemVer) descrito en la sección~\ref{sec:versionado} (pág.~\pageref{sec:versionado}).
\end{NotaTecnica}

\subsection{Catálogo de \textit{endpoints} por dominio}
\label{subsec:catalogo_endpoints}

A nivel de método, las anotaciones \comando{@GetMapping}, \comando{@PostMapping},
\comando{@PutMapping}, \comando{@PatchMapping} y \comando{@DeleteMapping} definen el
verbo HTTP y el sub-recurso. La tabla~\ref{tab:endpoints} detalla los \textit{endpoints}
de los dominios de mayor complejidad, extraídos del código fuente.

\vspace{0.2cm}
\noindent
\renewcommand{\arraystretch}{1.3}
\fontsize{8}{9.6}\selectfont\sffamily
\begin{tabularx}{\textwidth}{|C{1.4cm}|L{5.0cm}|Y|}
\hline
\rowcolor{headerTabla}
\sffamily\textbf{Verbo} & \sffamily\textbf{Ruta} & \sffamily\textbf{Operación} \\
\hline
POST & \pathbreak{/api/auth/register} & Registro con rol (\texttt{PROFESSIONAL}\,/\,\texttt{RECRUITER}). \\ \hline
\rowcolor{grisTecnico}
POST & \pathbreak{/api/auth/login} & Autenticación con credenciales; devuelve token y rol. \\ \hline
POST & \pathbreak{/api/auth/oauth/sync} & Sincroniza el perfil tras un \textit{login} OAuth de Supabase. \\ \hline
\rowcolor{grisTecnico}
GET & \pathbreak{/api/v1/cv-documents} & Lista los CV del usuario autenticado. \\ \hline
POST & \pathbreak{/api/v1/cv-documents} & Crea un nuevo documento de CV. \\ \hline
\rowcolor{grisTecnico}
PUT & \pathbreak{/api/v1/cv-documents/\{id\}} & Actualiza el contenido LaTeX de un CV. \\ \hline
DELETE & \pathbreak{/api/v1/cv-documents/\{id\}} & Elimina (lógicamente) un CV. \\ \hline
\rowcolor{grisTecnico}
POST & \pathbreak{/api/v1/cv-documents/ai-assist} & Asistencia IA sobre el CV (Gemini). \\ \hline
POST & \pathbreak{/api/v1/cv-documents/compile-latex} & Compila LaTeX a PDF (\texttt{application/pdf}). \\ \hline
\rowcolor{grisTecnico}
POST & \pathbreak{/api/v1/recruiter/ai-insights} & Genera \textit{insights} de talento (4 modos IA). \\ \hline
POST & \pathbreak{/api/v1/messages} & Envía un mensaje de chat (idempotente). \\ \hline
\rowcolor{grisTecnico}
POST & \pathbreak{/api/uploads} & Sube un archivo a Supabase Storage. \\ \hline
GET & \pathbreak{/api/admin/metrics/overview} & Métricas agregadas para el panel admin. \\ \hline
\rowcolor{grisTecnico}
GET & \pathbreak{/api/admin/metrics/growth} & Serie temporal de crecimiento de usuarios. \\ \hline
\end{tabularx}
\label{tab:endpoints}

\subsection{Envoltura de respuesta uniforme: \texorpdfstring{\texttt{ApiResponse<T>}}{ApiResponse}}
\label{subsec:apiresponse}

Todos los \textit{endpoints} —salvo la descarga binaria de PDF y la exportación de
cuenta— devuelven el mismo envoltorio genérico \comando{ApiResponse<T>}, un
\texttt{record} de Java que homogeneíza la forma de las respuestas y simplifica
drásticamente el consumo desde el frontend (que así trata todas las respuestas con un
único contrato, como se detalla en el capítulo~\ref{ch:frontend}, pág.~\pageref{ch:frontend}).

\begin{lstlisting}[language=Java, caption={Envoltorio de respuesta uniforme (\texttt{shared/ApiResponse.java}).}, label={lst:apiresponse}]
@JsonInclude(JsonInclude.Include.NON_NULL)
public record ApiResponse<T>(
    boolean success,      // true cuando la operacion tuvo exito
    int status,           // codigo HTTP reflejado en el cuerpo
    String message,       // mensaje legible para humanos
    T data,               // payload tipado; null en errores
    List<String> errors   // errores de validacion/negocio; null en exito
) {
    public static <T> ApiResponse<T> ok(T data, String message) { ... }
    public static <T> ApiResponse<T> created(T data, String message) { ... }
    public static <T> ApiResponse<T> badRequest(String m, List<String> e) { ... }
    // unauthorized(401), forbidden(403), conflict(409), internalError(500)
}
\end{lstlisting}

La anotación \comando{@JsonInclude(NON\_NULL)} omite los campos nulos en la
serialización: una respuesta de éxito no incluye el campo \texttt{errors} y una de error
no incluye \texttt{data}. El \textit{listado}~\ref{lst:json-ejemplo} contrasta una
respuesta de éxito y una de error de validación para el mismo \textit{endpoint} de
registro. La tabla~\ref{tab:apiresponse-campos} documenta cada campo.

\clearpage
\begin{lstlisting}[language=Java, caption={Cuerpos JSON de éxito (201) y de error de validación (400).}, label={lst:json-ejemplo}]
// 201 Created - registro exitoso (sin campo "errors")
{ "success": true, "status": 201, "message": "Cuenta creada",
  "data": { "token": "eyJ...", "tokenType": "Bearer", "role": "PROFESSIONAL" } }

// 400 Bad Request - validacion fallida (sin campo "data")
{ "success": false, "status": 400, "message": "Datos invalidos",
  "errors": ["email: formato invalido", "password: minimo 8 caracteres"] }
\end{lstlisting}

\vspace{0.2cm}
\noindent
\renewcommand{\arraystretch}{1.3}
\fontsize{8.5}{10.2}\selectfont\sffamily
\begin{tabularx}{\textwidth}{|L{2.2cm}|C{2.4cm}|Y|}
\hline
\rowcolor{headerTabla}
\sffamily\textbf{Campo} & \sffamily\textbf{Tipo} & \sffamily\textbf{Semántica} \\
\hline
\texttt{success} & \texttt{boolean} & Indicador binario de éxito; permite al cliente bifurcar sin inspeccionar el código HTTP. \\ \hline
\rowcolor{grisTecnico}
\texttt{status} & \texttt{int} & Código HTTP reflejado también en el cuerpo (redundancia útil para logs de cliente). \\ \hline
\texttt{message} & \texttt{String} & Mensaje legible, ya localizado en español cuando procede. \\ \hline
\rowcolor{grisTecnico}
\texttt{data} & \texttt{T} (genérico) & Carga útil tipada; \texttt{null} en respuestas de error. \\ \hline
\texttt{errors} & \texttt{List<String>} & Lista de errores de validación o negocio; \texttt{null} en éxito. \\ \hline
\end{tabularx}
\label{tab:apiresponse-campos}

\subsection{Autogeneración de especificaciones (springdoc-openapi)}
\label{subsec:openapi}

La documentación de la API \textbf{no se mantiene a mano}: se genera automáticamente a
partir de las anotaciones del código mediante \textbf{springdoc-openapi 2.7}. Las
anotaciones \comando{@Operation}, \comando{@Tag} y \comando{@Schema}, presentes en los
controladores y DTOs reales, alimentan una especificación OpenAPI~3 navegable y siempre
sincronizada con la implementación.

\vspace{0.2cm}
\noindent
\renewcommand{\arraystretch}{1.3}
\fontsize{9}{10.8}\selectfont\sffamily
\begin{tabularx}{\textwidth}{|L{4.0cm}|Y|}
\hline
\rowcolor{headerTabla}
\sffamily\textbf{Recurso} & \sffamily\textbf{Descripción} \\
\hline
\pathbreak{/swagger-ui.html} & Interfaz interactiva Swagger UI para explorar y probar los \textit{endpoints}. \\
\hline
\rowcolor{grisTecnico}
\pathbreak{/v3/api-docs} & Especificación OpenAPI~3 en formato JSON, consumible por herramientas externas. \\
\hline
\end{tabularx}

\begin{AvisoNota}
Ambas rutas figuran en la lista \texttt{PUBLIC\_PATHS} de \texttt{SecurityConfig}
(\pathbreak{/swagger-ui/**} y \pathbreak{/v3/api-docs/**}), por lo que la documentación
es accesible sin autenticación en los entornos donde se habilite. El esquema de seguridad
\texttt{bearerAuth} queda declarado para que Swagger UI permita adjuntar un JWT en las
pruebas interactivas.
\end{AvisoNota}

\subsection{Convenciones de códigos de estado HTTP}
\label{subsec:http_status}

La API aplica los códigos de estado HTTP de forma \textbf{semántica y consistente}, lo que
permite al cliente reaccionar de manera uniforme (véase el interceptor de respuesta del
capítulo~\ref{ch:frontend}, sección~\ref{subsec:interceptor_resp},
pág.~\pageref{subsec:interceptor_resp}). La tabla~\ref{tab:http-status} documenta el
contrato de estados.

\clearpage
\vspace{0.2cm}
\noindent
\renewcommand{\arraystretch}{1.3}
\fontsize{8.5}{10.2}\selectfont\sffamily
\begin{tabularx}{\textwidth}{|C{1.6cm}|L{3.0cm}|Y|}
\hline
\rowcolor{headerTabla}
\sffamily\textbf{Código} & \sffamily\textbf{Significado} & \sffamily\textbf{Cuándo se emite} \\
\hline
200 & OK & Operación de lectura o actualización exitosa. \\ \hline
\rowcolor{grisTecnico}
201 & Created & Creación exitosa de un recurso (\texttt{ApiResponse.created}). \\ \hline
400 & Bad Request & Validación fallida; incluye lista de \texttt{errors} por campo. \\ \hline
\rowcolor{grisTecnico}
401 & Unauthorized & Token ausente, inválido o caducado. \\ \hline
403 & Forbidden & Rol insuficiente o acceso a recurso ajeno. \\ \hline
\rowcolor{grisTecnico}
409 & Conflict & Violación de unicidad (p.~ej., correo ya registrado). \\ \hline
422 & Unprocessable & Compilación LaTeX fallida en el CV Studio. \\ \hline
\rowcolor{grisTecnico}
429 & Too Many Requests & Límite de uso de IA alcanzado (Gemini). \\ \hline
503 & Service Unavailable & Dependencia externa (IA) no disponible. \\ \hline
\end{tabularx}
\label{tab:http-status}

\subsection{Superficie administrativa y métricas}
\label{subsec:admin}

La familia \texttt{Admin*Controller} expone una superficie de \textbf{observabilidad y
moderación} accesible solo al rol \texttt{ROLE\_ADMIN}. El \comando{AdminMetricsController}
agrega indicadores de negocio mediante consultas a la base de datos, resumidos en la
tabla~\ref{tab:admin-metrics}.

\vspace{0.2cm}
\noindent
\renewcommand{\arraystretch}{1.3}
\fontsize{8}{9.6}\selectfont\sffamily
\begin{tabularx}{\textwidth}{|L{5.2cm}|Y|}
\hline
\rowcolor{headerTabla}
\sffamily\textbf{Endpoint} & \sffamily\textbf{Métrica} \\
\hline
\pathbreak{GET /api/admin/metrics/overview} & Resumen agregado (usuarios, perfiles, actividad). \\ \hline
\rowcolor{grisTecnico}
\pathbreak{GET /api/admin/metrics/growth} & Serie temporal de crecimiento de usuarios. \\ \hline
\pathbreak{GET /api/admin/metrics/roles} & Distribución de usuarios por rol. \\ \hline
\rowcolor{grisTecnico}
\pathbreak{GET /api/admin/metrics/top-skills} & Ranking de habilidades más frecuentes. \\ \hline
\pathbreak{GET /api/admin/metrics/health} & Indicadores de salud del sistema. \\ \hline
\rowcolor{grisTecnico}
\pathbreak{GET /api/admin/logs} & Registro de eventos para auditoría. \\ \hline
\end{tabularx}
\label{tab:admin-metrics}

Las demás familias administrativas atienden ámbitos específicos:
\pathbreak{/api/admin/moderation} (moderación de contenido reportado),
\pathbreak{/api/admin/profiles} (gestión de perfiles), \pathbreak{/api/admin/skills}
(catálogo curado de \textit{hard}/\textit{soft skills}) y \pathbreak{/api/admin/email}
(envíos administrativos). Todas comparten la protección por rol descrita en la
sección~\ref{sec:excepciones} (pág.~\pageref{sec:excepciones}).

\subsection{Referencia de \textit{endpoints} de los dominios de contenido}
\label{subsec:endpoints_contenido}

Para completar la trazabilidad, la tabla~\ref{tab:endpoints-contenido} detalla los
\textit{endpoints} reales de los dominios de contenido profesional —proyectos, habilidades,
experiencia, formación y portafolio—, extraídos de las anotaciones de método de sus
respectivos controladores.

\clearpage
\vspace{0.2cm}
\noindent
\renewcommand{\arraystretch}{1.25}
\fontsize{7.6}{9.2}\selectfont\sffamily
\begin{tabularx}{\textwidth}{|C{1.3cm}|L{5.6cm}|Y|}
\hline
\rowcolor{headerTabla}
\sffamily\textbf{Verbo} & \sffamily\textbf{Ruta} & \sffamily\textbf{Operación} \\
\hline
GET & \pathbreak{/api/projects/profile/\{profileId\}} & Proyectos de un perfil. \\ \hline
\rowcolor{grisTecnico}
POST & \pathbreak{/api/projects} & Crea un proyecto. \\ \hline
DELETE & \pathbreak{/api/projects/\{projectId\}} & Elimina un proyecto. \\ \hline
\rowcolor{grisTecnico}
GET & \pathbreak{/api/profiles/\{profileId\}/skills/hard} & \textit{Hard skills} de un perfil. \\ \hline
PATCH & \pathbreak{/api/skills/hard/\{skillId\}/top} & Marca una \textit{skill} como destacada. \\ \hline
\rowcolor{grisTecnico}
PATCH & \pathbreak{/api/profiles/\{profileId\}/skills/hard/top-order} & Reordena las \textit{skills} destacadas. \\ \hline
PATCH & \pathbreak{/api/v1/profile/basic} & Actualiza datos básicos del perfil. \\ \hline
\rowcolor{grisTecnico}
PATCH & \pathbreak{/api/v1/profile/bio} & Actualiza la biografía. \\ \hline
POST & \pathbreak{/api/v1/work-experiences} & Crea una experiencia (con lat/lng). \\ \hline
\rowcolor{grisTecnico}
PUT & \pathbreak{/api/v1/portfolio/settings} & Guarda los ajustes del portafolio. \\ \hline
GET & \pathbreak{/api/v1/portfolio/export} & Exporta el portafolio. \\ \hline
\rowcolor{grisTecnico}
POST & \pathbreak{/api/v1/portfolio/curriculum/compile} & Compila el currículum del portafolio a PDF. \\ \hline
\end{tabularx}
\label{tab:endpoints-contenido}

\begin{NotaTecnica}
Las rutas de borrado de experiencia y formación incluyen el \texttt{profileId} en la URL
(\pathbreak{/\{id\}/profile/\{profileId\}}) como verificación redundante de propiedad: el
servidor comprueba que el recurso pertenezca a ese perfil antes de eliminarlo, una defensa
adicional sobre la ya provista por las políticas RLS del motor
(capítulo~\ref{ch:basedatos}, pág.~\pageref{ch:basedatos}).
\end{NotaTecnica}
        % 4.1 Arquitectura de la API REST e interfaces
% =============================================================================
% 4.2 DIAGRAMAS DE CLASES POR DOMINIOS CRÍTICOS
% =============================================================================
\section{Diagramas de Clases por Dominios Críticos}
\label{sec:clases}

Esta sección documenta la \textbf{estructura estática} del backend: cómo se organizan las
clases dentro de cada dominio y cómo fluyen los datos entre ellas. El patrón es uniforme
y deliberado, lo que reduce la carga cognitiva al navegar el código: \textit{todo dominio
se lee de la misma forma}.

\subsection{El patrón de capas por \textit{feature}}
\label{subsec:patron_capas}

Cada dominio implementa tres responsabilidades bien separadas, según el principio de
responsabilidad única (SRP):

\begin{itemize}[noitemsep]
    \item \comando{@RestController}: frontera HTTP. Recibe DTOs, valida y delega; nunca
    contiene lógica de negocio ni acceso a datos.
    \item \comando{@Service}: lógica de negocio y límite transaccional
    (\comando{@Transactional}). Orquesta repositorios y servicios externos.
    \item \comando{Repository}: acceso a datos tipado con jOOQ a través de un único
    \comando{DSLContext} inyectado por Spring.
\end{itemize}

La figura~\ref{fig:clases-capas} ilustra el patrón sobre tres dominios críticos
—\texttt{profile}, \texttt{portfolio} y \texttt{cvstudio}— y evidencia la dependencia
descendente Controller~$\rightarrow$~Service~$\rightarrow$~Repository.

El backend organiza el código en \textbf{paquetes por dominio} bajo
\pathbreak{com.ethoshub}, cada uno autocontenido. La figura~\ref{fig:paquetes-backend}
ofrece un mapa de los paquetes reales, agrupados por su rol funcional.

\vspace{0.3cm}
\begin{figure}[h!]
\centering
\begin{tikzpicture}[
    pkg/.style={rectangle, rounded corners=2pt, draw=c4Sistema!70!black, fill=c4Componente!28,
        font=\sffamily\scriptsize, align=center, minimum width=2.5cm, minimum height=0.7cm, inner sep=2pt},
    grp/.style={rectangle, rounded corners=4pt, draw=grisISO!60, dashed, inner sep=7pt},
    lbl/.style={font=\sffamily\scriptsize\bfseries, text=colorPrimario}]
    % Identidad
    \node[pkg] (auth) at (0,3) {auth};
    \node[pkg] (sec) at (2.6,3) {security};
    \node[pkg] (prof) at (5.2,3) {profile};
    \node[pkg] (pref) at (7.8,3) {preferences};
    \node[lbl] at (3.9,3.7) {Identidad y cuentas};
    \begin{scope}[on background layer]
      \node[grp, fit=(auth)(pref), draw=c4Persona!60] {};
    \end{scope}
    % Contenido
    \node[pkg] (proj) at (0,1.4) {project};
    \node[pkg] (skill) at (2.6,1.4) {skills};
    \node[pkg] (exp) at (5.2,1.4) {experience};
    \node[pkg] (edu) at (7.8,1.4) {education};
    \node[pkg] (port) at (1.3,0.4) {portfolio};
    \node[pkg] (cv) at (3.9,0.4) {cvstudio};
    \node[lbl] at (3.9,2.1) {Contenido profesional};
    % Interacciones
    \node[pkg] (chat) at (0,-1.1) {chat};
    \node[pkg] (notif) at (2.6,-1.1) {notifications};
    \node[pkg] (conn) at (5.2,-1.1) {connections};
    \node[pkg] (recr) at (7.8,-1.1) {recruiter (+ ai)};
    \node[lbl] at (3.9,-0.4) {Interacciones};
    % Transversal
    \node[pkg, fill=colorAcento!20] (shared) at (1.3,-2.4) {shared};
    \node[pkg, fill=colorAcento!20] (config) at (3.9,-2.4) {config};
    \node[pkg, fill=colorAcento!20] (admin) at (6.5,-2.4) {admin};
    \node[lbl] at (3.9,-1.8) {Transversal};
\end{tikzpicture}
\caption{Mapa de paquetes del backend bajo \texttt{com.ethoshub}, agrupados por rol.}
\label{fig:paquetes-backend}
\end{figure}

\vspace{0.3cm}
\begin{figure}[h!]
\centering
\begin{tikzpicture}[
    cls/.style={rectangle, draw=c4Sistema!70!black, fill=c4Componente!25, font=\sffamily\scriptsize,
        align=center, minimum width=2.7cm, minimum height=0.8cm, inner sep=3pt},
    ctrl/.style={cls, fill=c4Componente!50},
    rel/.style={-{Stealth[length=2mm]}, semithick, draw=grisISO},
    node distance=0.7cm and 0.5cm
]
    \node[ctrl] (pc) at (0,3) {\texttt{ProfileController}};
    \node[cls] (ps) at (0,1.8) {\texttt{ProfileService}};
    \node[cls] (pr) at (0,0.6) {\texttt{ProfileRepository}};
    \node[ctrl] (bc) at (3.6,3) {\texttt{PortfolioSettings}\\\texttt{Controller}};
    \node[cls] (bs) at (3.6,1.8) {\texttt{PortfolioSettings}\\\texttt{Service}};
    \node[cls] (br) at (3.6,0.6) {\texttt{PortfolioRepository}};
    \node[ctrl] (cc) at (7.2,3) {\texttt{CvDocument}\\\texttt{Controller}};
    \node[cls] (cs) at (7.2,1.8) {\texttt{CvDocumentService}\\\texttt{+ LatexCompileSvc}};
    \node[cls] (cr) at (7.2,0.6) {\texttt{CvDocumentRepository}};

    \node[cls, fill=c4Datos!25, draw=c4Datos!70!black, minimum width=8.5cm] (dsl) at (3.6,-0.7)
        {jOOQ \texttt{DSLContext} $\rightarrow$ PostgreSQL \texttt{core}};

    \foreach \a/\b in {pc/ps, ps/pr, bc/bs, bs/br, cc/cs, cs/cr} \draw[rel] (\a) -- (\b);
    \draw[rel] (pr) -- (dsl); \draw[rel] (br) -- (dsl); \draw[rel] (cr) -- (dsl);
\end{tikzpicture}
\caption{Patrón de capas por \textit{feature} en los dominios \texttt{profile}, \texttt{portfolio} y \texttt{cvstudio}.}
\label{fig:clases-capas}
\end{figure}

El dominio \texttt{cvstudio} es el de mayor riqueza estructural: su servicio coordina no
solo el repositorio, sino dos colaboradores especializados —\comando{GeminiAiService}
(asistencia de IA) y \comando{LatexCompileService} (compilación a PDF)—, ambos
documentados en la sección~\ref{sec:secuencias} (pág.~\pageref{sec:secuencias}).

\subsection{Modelado de DTOs y registros tipados de jOOQ}
\label{subsec:dtos}

Los \textbf{objetos de transferencia de datos} (DTOs) se implementan como
\comando{record} de Java inmutables, con validación declarativa mediante \textbf{Jakarta
Bean Validation}. Se distinguen dos familias:

\begin{itemize}[noitemsep]
    \item \textbf{DTOs de petición} (\texttt{*Request}): anotados con restricciones
    (\comando{@NotBlank}, \comando{@Email}, \comando{@Size}) que Spring evalúa
    automáticamente al deserializar el cuerpo JSON.
    \item \textbf{DTOs de respuesta} (\texttt{*Response} / \texttt{*Dto}): proyecciones de
    solo lectura que el servicio construye a partir de los registros tipados de jOOQ.
\end{itemize}

Entre ambos median los \textbf{registros generados por jOOQ} a partir del esquema real de
la base de datos: clases como \texttt{ProfilesBasicRecord} o \texttt{CvDocumentsRecord}
que ofrecen acceso tipado y verificado en compilación a cada columna, eliminando los
errores de SQL por cadenas mágicas.

La figura~\ref{fig:dto-flujo} muestra el viaje de los datos desde el DTO de petición
validado hasta el DTO de respuesta serializado.

\vspace{0.3cm}
\begin{figure}[h!]
\centering
\begin{tikzpicture}[node distance=0.6cm and 1.3cm,
    box/.style={rectangle, rounded corners=2pt, draw, font=\sffamily\scriptsize, align=center,
        minimum width=2.6cm, minimum height=0.85cm, inner sep=3pt},
    rel/.style={-{Stealth[length=2mm]}, semithick}]
    \node[box, fill=c4Componente!40] (req) {\texttt{RegisterRequest}\\\scriptsize @NotBlank @Email @Size};
    \node[box, fill=colorAcento!30, right=of req] (val) {Bean Validation\\\scriptsize (Spring)};
    \node[box, fill=c4Componente!40, right=of val] (svc) {\texttt{AuthService}\\\scriptsize lógica};
    \node[box, fill=c4Datos!25, below=of svc] (db) {jOOQ\\\texttt{ProfilesBasicRecord}};
    \node[box, fill=c4Componente!40, below=of req] (resp) {\texttt{AuthResponse}\\\scriptsize token + rol};

    \draw[rel] (req) -- node[c4etiqueta, above]{JSON} (val);
    \draw[rel] (val) -- node[c4etiqueta, above]{válido} (svc);
    \draw[rel] (svc) -- (db);
    \draw[rel] (svc) to[bend left=10] (resp);
    \draw[rel, dashed] (val) to[bend right=30] node[c4etiqueta, left]{400 +\\errores} (resp);
\end{tikzpicture}
\caption{Flujo de un DTO de petición validado hacia el DTO de respuesta.}
\label{fig:dto-flujo}
\end{figure}

\begin{lstlisting}[language=Java, caption={DTO de registro con validación declarativa (\texttt{auth/dto/RegisterRequest.java}).}, label={lst:registerrequest}]
public record RegisterRequest(
    @NotBlank @Email                 String email,
    @NotBlank @Size(min = 8)         String password,
    @NotNull                         UserRole role,   // PROFESSIONAL | RECRUITER
    String firstName, String lastName,
    String phoneCode, String phoneNumber
) { }
\end{lstlisting}

Cuando la validación falla, \comando{MethodArgumentNotValidException} es interceptada por
el manejador global (sección~\ref{sec:excepciones}, pág.~\pageref{sec:excepciones}) y
traducida a una respuesta \texttt{400} con la lista de errores por campo, sin que el
controlador escriba una sola línea de código de validación manual.

\begin{NotaTecnica}
La elección de \texttt{record} inmutables para los DTOs no es estética: garantiza que un
objeto de petición no pueda mutar tras su validación, eliminando una clase entera de
errores (\textit{time-of-check to time-of-use}). La inmutabilidad también facilita el
razonamiento concurrente, relevante para el modelo \textit{stateless} del servidor.
\end{NotaTecnica}

\subsection{Resumen de DTOs por dominio crítico}
\label{subsec:dtos_resumen}

La tabla~\ref{tab:dtos} cataloga los DTOs más representativos y su rol en el contrato de
la API.

\vspace{0.2cm}
\noindent
\renewcommand{\arraystretch}{1.3}
\fontsize{8.5}{10.2}\selectfont\sffamily
\begin{tabularx}{\textwidth}{|L{4.2cm}|C{1.8cm}|Y|}
\hline
\rowcolor{headerTabla}
\sffamily\textbf{DTO} & \sffamily\textbf{Familia} & \sffamily\textbf{Uso} \\
\hline
\texttt{RegisterRequest} & Petición & Alta de usuario con rol y datos de contacto. \\ \hline
\rowcolor{grisTecnico}
\texttt{LoginRequest} & Petición & Credenciales para \textit{login} con email/contraseña. \\ \hline
\texttt{AuthResponse} & Respuesta & Token, tipo de token, expiración y rol resuelto. \\ \hline
\rowcolor{grisTecnico}
\texttt{CvDocumentRequest} & Petición & Título y código LaTeX del CV a crear/actualizar. \\ \hline
\texttt{AiAssistRequest} & Petición & Prompt, contenido, modo e historial para Gemini. \\ \hline
\rowcolor{grisTecnico}
\texttt{TalentDto} & Respuesta & Proyección de un perfil profesional para reclutadores. \\ \hline
\texttt{ChatMessageRequest} & Petición & Mensaje, \texttt{chatId} y clave de idempotencia. \\ \hline
\end{tabularx}
\label{tab:dtos}

\subsection{Diagrama de clases UML del dominio \texttt{cvstudio}}
\label{subsec:uml_cvstudio}

La figura~\ref{fig:uml-cvstudio} ofrece una vista UML del dominio más rico del backend,
con sus atributos y operaciones públicas. Evidencia cómo \comando{CvDocumentService}
coordina tres colaboradores y traduce entre la entidad de dominio y los DTOs.

\clearpage
\vspace{0.3cm}
\begin{figure}[h!]
\centering
\begin{tikzpicture}[
    uml/.style={rectangle, draw=c4Sistema!70!black, fill=white, font=\sffamily\scriptsize,
        align=left, inner sep=4pt, minimum width=3.6cm},
    titulo/.style={font=\sffamily\scriptsize\bfseries, text=colorPrimario},
    rel/.style={-{Stealth[length=2mm]}, semithick, draw=grisISO},
    agg/.style={-{Diamond[open, length=3mm]}, semithick, draw=grisISO}]

    \node[uml] (ctrl) at (0,3.4) {%
        {\titulo CvDocumentController}\\\hrule height0.4pt\vspace{2pt}
        + list() : ApiResponse\\
        + create(req) : ApiResponse\\
        + aiAssist(req) : ApiResponse\\
        + compileLatex(req) : byte[]};
    \node[uml] (svc) at (0,0.6) {%
        {\titulo CvDocumentService}\\\hrule height0.4pt\vspace{2pt}
        - repo : CvDocumentRepository\\
        - gemini : GeminiAiService\\
        - latex : LatexCompileService\\\hrule height0.4pt\vspace{2pt}
        + assist(...) : String\\
        + compile(...) : byte[]};
    \node[uml, fill=c4Datos!12] (repo) at (-3.6,-2.4) {%
        {\titulo CvDocumentRepository}\\\hrule height0.4pt\vspace{2pt}
        - dsl : DSLContext\\
        + findByProfile(id)\\
        + upsert(rec)};
    \node[uml, fill=colorAcento!10] (gem) at (0,-2.4) {%
        {\titulo GeminiAiService}\\\hrule height0.4pt\vspace{2pt}
        + generate\\Assistance(...)};
    \node[uml, fill=colorAcento!10] (tex) at (3.6,-2.4) {%
        {\titulo LatexCompileService}\\\hrule height0.4pt\vspace{2pt}
        + compile(code,\\\ \ fotoUrl)};

    \draw[rel] (ctrl) -- (svc);
    \draw[agg] (svc) -- (repo);
    \draw[agg] (svc) -- (gem);
    \draw[agg] (svc) -- (tex);
\end{tikzpicture}
\caption{Diagrama de clases UML del dominio \texttt{cvstudio} (atributos y operaciones).}
\label{fig:uml-cvstudio}
\end{figure}

\subsection{Acceso a datos tipado con jOOQ}
\label{subsec:jooq}

Los repositorios no escriben SQL como cadenas de texto: usan el \textbf{DSL tipado} de
jOOQ, generado a partir del esquema real de la base de datos. El \textit{listado}
\ref{lst:jooq} muestra un repositorio representativo donde cada nombre de tabla y columna
está verificado en tiempo de compilación.

\begin{lstlisting}[language=Java, caption={Acceso tipado con jOOQ en un repositorio (patrón real).}, label={lst:jooq}]
@Repository
public class CvDocumentRepository {
    private final DSLContext dsl;
    public List<CvDocumentsRecord> findByProfile(UUID profileId) {
        return dsl.selectFrom(CV_DOCUMENTS)
                  .where(CV_DOCUMENTS.PROFILE_ID.eq(profileId))
                  .and(CV_DOCUMENTS.DELETED_AT.isNull())   // borrado logico
                  .orderBy(CV_DOCUMENTS.UPDATED_AT.desc())
                  .fetch();
    }
}
\end{lstlisting}

\clearpage
\begin{NotaTecnica}
La ventaja decisiva de jOOQ sobre el SQL textual o un ORM por reflexión es la
\textbf{seguridad en compilación}: si una columna se renombra en la base de datos y se
regeneran las clases, cualquier consulta que la referencie deja de compilar. Esto convierte
una clase entera de errores de tiempo de ejecución en errores detectados antes del
despliegue, en línea con la estrategia de verificación del capítulo~8.
\end{NotaTecnica}
          % 4.2 Diagramas de clases por dominios
% =============================================================================
% 4.3 DIAGRAMAS DE SECUENCIA DE OPERACIONES CLAVE
% =============================================================================
\section{Diagramas de Secuencia de Operaciones Clave}
\label{sec:secuencias}

Esta sección documenta el \textbf{comportamiento dinámico} del backend mediante diagramas
de secuencia de las tres operaciones de mayor valor y complejidad: la autenticación
federada con Supabase, el ciclo de asistencia IA y compilación de currículums, y la
generación de \textit{insights} de reclutamiento.

\subsection{Autenticación federada y validación de JWT}
\label{subsec:seq_auth}

El backend \textbf{nunca emite tokens}: delega la identidad en Supabase Auth y se limita
a \textbf{verificar} cada JWT contra el JWKS (\textit{JSON Web Key Set}) del proveedor.
La figura~\ref{fig:seq-auth} detalla la secuencia para una petición protegida.

\vspace{0.3cm}
\begin{figure}[h!]
\centering
\begin{tikzpicture}[
    lifeline/.style={font=\sffamily\scriptsize\bfseries, align=center},
    msg/.style={-{Stealth[length=2mm]}, semithick},
    ret/.style={-{Stealth[length=2mm]}, semithick, dashed},
    note/.style={font=\sffamily\scriptsize\itshape, text=grisISO}
]
    \def\xspa{0} \def\xfil{3.2} \def\xjwks{6.4} \def\xctrl{9.0}
    \node[lifeline] (spa) at (\xspa,0) {SPA\\(Axios)};
    \node[lifeline] (fil) at (\xfil,0) {Spring Security\\(Resource Server)};
    \node[lifeline] (jwks) at (\xjwks,0) {Supabase\\JWKS};
    \node[lifeline] (ctrl) at (\xctrl,0) {Controller\\+ Service};
    \foreach \x in {\xspa,\xfil,\xjwks,\xctrl} \draw[grisISO!50, thick] (\x,-0.4) -- (\x,-6.2);

    \draw[msg] (\xspa,-0.9) -- node[note, above]{\scriptsize GET /api/... + Bearer JWT} (\xfil,-0.9);
    \draw[msg] (\xfil,-1.7) -- node[note, above]{\scriptsize obtener claves (1ª vez)} (\xjwks,-1.7);
    \draw[ret] (\xjwks,-2.3) -- node[note, above]{\scriptsize claves públicas (cacheadas)} (\xfil,-2.3);
    \node[note, align=left] at (\xfil+0.1,-3.0) {\scriptsize verifica firma (RS256/HS256),\\\scriptsize exp y claims};
    \node[note, align=left] at (\xfil+0.1,-3.8) {\scriptsize AdminEmailPolicy.resolveRole()};
    \draw[msg] (\xfil,-4.4) -- node[note, above]{\scriptsize Authentication + authority} (\xctrl,-4.4);
    \draw[ret] (\xctrl,-5.1) -- node[note, above]{\scriptsize ApiResponse<T> (JSON)} (\xspa,-5.1);
    \node[note, align=left] at (\xfil+0.2,-5.8) {\scriptsize si falla: 401 (entryPoint) o 403 (accessDenied)};
\end{tikzpicture}
\caption{Secuencia de validación de un JWT de Supabase en una petición protegida.}
\label{fig:seq-auth}
\end{figure}

El paso de \textbf{resolución de rol} es crítico: \comando{AdminEmailPolicy.resolveRole()}
deriva la autoridad \\ (\texttt{ROLE\_PROFESSIONAL}, \texttt{ROLE\_RECRUITER} o
\texttt{ROLE\_ADMIN}) a partir de los \textit{claims} del token y degrada cualquier
intento de \texttt{role=admin} que no provenga de un correo administrativo autorizado.
Las claves se pre-descargan al arranque (\comando{warmupJwks()}) para que la primera
petición no pague la latencia de la primera resolución JWKS.

\clearpage
\subsection{Recuperación de contraseña mediante OTP por correo}
\label{subsec:seq_otp}

El reinicio de contraseña no expone enlaces con tokens en la URL: usa un \textbf{código de
un solo uso} (OTP) enviado por correo, gestionado por \comando{ForgotPasswordController} y
\comando{ForgotPasswordService} en tres pasos (\texttt{/request}, \texttt{/verify},
\texttt{/reset}). La figura~\ref{fig:seq-otp} detalla el flujo, que persiste el código en
\texttt{core.verification\_codes} (tabla blindada por RLS a \texttt{service\_role}).

\vspace{0.3cm}
\begin{figure}[h!]
\centering
\begin{tikzpicture}[
    lifeline/.style={font=\sffamily\scriptsize\bfseries, align=center},
    msg/.style={-{Stealth[length=2mm]}, semithick},
    ret/.style={-{Stealth[length=2mm]}, semithick, dashed},
    note/.style={font=\sffamily\scriptsize\itshape, text=grisISO}
]
    \def\xui{0} \def\xc{3.4} \def\xs{6.6} \def\xm{9.4}
    \node[lifeline] (ui) at (\xui,0) {SPA};
    \node[lifeline] (c) at (\xc,0) {ForgotPassword\\Controller};
    \node[lifeline] (s) at (\xs,0) {Service +\\verification\_codes};
    \node[lifeline] (m) at (\xm,0) {EmailService\\(SMTP)};
    \foreach \x in {\xui,\xc,\xs,\xm} \draw[grisISO!50, thick] (\x,-0.4) -- (\x,-5.8);

    \draw[msg] (\xui,-0.9) -- node[note, above]{\scriptsize POST /request} (\xc,-0.9);
    \draw[msg] (\xc,-1.5) -- node[note, above]{\scriptsize genera OTP} (\xs,-1.5);
    \draw[msg] (\xs,-2.1) -- node[note, above]{\scriptsize envía código} (\xm,-2.1);
    \draw[ret] (\xc,-2.7) -- node[note, above]{\scriptsize 200 (siempre)} (\xui,-2.7);
    \draw[msg] (\xui,-3.4) -- node[note, above]{\scriptsize POST /verify (código)} (\xc,-3.4);
    \draw[msg] (\xc,-4.0) -- node[note, above]{\scriptsize valida + expira} (\xs,-4.0);
    \draw[msg] (\xui,-4.7) -- node[note, above]{\scriptsize POST /reset (nueva clave)} (\xc,-4.7);
    \draw[ret] (\xc,-5.3) -- node[note, above]{\scriptsize ApiResponse OK} (\xui,-5.3);
\end{tikzpicture}
\caption{Secuencia de recuperación de contraseña con código OTP por correo.}
\label{fig:seq-otp}
\end{figure}

\begin{AlertaSeguridad}
El \textit{endpoint} \pathbreak{/request} devuelve \texttt{200} \textbf{siempre}, exista o
no la cuenta, para no revelar qué correos están registrados (mitigación de enumeración de
usuarios). El código OTP tiene caducidad y un único uso: \comando{record\_login\_failure}
y la lógica de bloqueo (cabecera \texttt{X-Blocked-Until}) limitan los intentos de fuerza
bruta.
\end{AlertaSeguridad}

\subsection{Subida de archivos a Supabase Storage}
\label{subsec:seq_upload}

La subida de imágenes (fotos de perfil, medios de proyecto) la gestiona
\comando{FileUploadController} (\pathbreak{POST /api/uploads}) delegando en
\comando{SupabaseStorageService}, que envuelve la API REST de Storage sobre el \textit{bucket}
\texttt{profile-assets} con semántica de \textit{upsert}. El servidor actúa como
\textbf{intermediario de confianza}: el cliente nunca recibe la \texttt{service\_role key};
es el backend quien firma la petición a Storage. Spring limita el tamaño a \textbf{20\,MB}
por archivo (\texttt{max-file-size}) y \textbf{25\,MB} por petición.

\subsection{Generación y exportación de currículums en el CV Studio}
\label{subsec:seq_cv}

El CV Studio es la funcionalidad de IA más rica del sistema. Combina dos operaciones
asíncronas: la \textbf{asistencia generativa} (Gemini reescribe o sugiere texto) y la
\textbf{compilación a PDF} mediante un compilador LaTeX externo. La
figura~\ref{fig:seq-cv} muestra ambos ciclos.

\vspace{0.3cm}
\begin{figure}[h!]
\centering
\begin{tikzpicture}[
    lifeline/.style={font=\sffamily\scriptsize\bfseries, align=center},
    msg/.style={-{Stealth[length=2mm]}, semithick},
    ret/.style={-{Stealth[length=2mm]}, semithick, dashed},
    note/.style={font=\sffamily\scriptsize\itshape, text=grisISO}
]
    \def\xui{0} \def\xc{3.4} \def\xs{6.6} \def\xg{9.4}
    \node[lifeline] (ui) at (\xui,0) {CV Studio\\(Monaco)};
    \node[lifeline] (c) at (\xc,0) {CvDocument\\Controller};
    \node[lifeline] (s) at (\xs,0) {Gemini /\\LatexCompile Svc};
    \node[lifeline] (g) at (\xg,0) {Gemini API /\\Compilador LaTeX};
    \foreach \x in {\xui,\xc,\xs,\xg} \draw[grisISO!50, thick] (\x,-0.4) -- (\x,-6.0);

    \draw[msg] (\xui,-0.9) -- node[note, above]{\scriptsize POST /ai-assist} (\xc,-0.9);
    \draw[msg] (\xc,-1.5) -- (\xs,-1.5);
    \draw[msg] (\xs,-2.1) -- node[note, above]{\scriptsize prompt + historial} (\xg,-2.1);
    \draw[ret] (\xg,-2.7) -- node[note, above]{\scriptsize sugerencia} (\xs,-2.7);
    \draw[ret] (\xs,-3.3) -- node[note, above]{\scriptsize ApiResponse} (\xui,-3.3);
    \draw[msg] (\xui,-4.0) -- node[note, above]{\scriptsize POST /compile-latex} (\xc,-4.0);
    \draw[msg] (\xc,-4.6) -- node[note, above]{\scriptsize compile(code, fotoUrl)} (\xs,-4.6);
    \draw[msg] (\xs,-5.1) -- node[note, above]{\scriptsize proceso externo} (\xg,-5.1);
    \draw[ret] (\xs,-5.7) -- node[note, above]{\scriptsize PDF (binario, 200) / 422 si falla} (\xui,-5.7);
\end{tikzpicture}
\caption{Secuencia de asistencia IA y compilación de CV a PDF en el CV Studio.}
\label{fig:seq-cv}
\end{figure}

\paragraph{Compilación LaTeX en proceso aislado.} El servicio
\comando{LatexCompileService} no compila en el hilo de la JVM: crea un \textbf{directorio
temporal único} (\texttt{tectonic-\{uuid\}}), \textbf{sanea} el código fuente
(\comando{UnicodeLatexSanitizer}) para neutralizar caracteres Unicode problemáticos,
inyecta la foto de perfil mediante el marcador \comando{\textbackslash ethoshubFotoPerfil}
y lanza el binario del compilador como \textbf{proceso externo} con un \textit{timeout}
de \textbf{120 segundos}. El binario se resuelve dinámicamente probando rutas candidatas
(\pathbreak{/usr/bin/pdflatex}, \pathbreak{/usr/local/bin/pdflatex}). Si el proceso falla
o expira, se devuelve \texttt{422} en lugar de propagar una excepción genérica.

\begin{AlertaSeguridad}
Compilar LaTeX arbitrario de usuario es un vector de riesgo (ejecución de \textit{shell}
vía \texttt{\textbackslash write18}). La mitigación combina tres barreras: saneamiento
previo del código, ejecución en directorio temporal efímero que se elimina al terminar
(\texttt{Files.walk(...).sorted(reverseOrder)}), y \textit{timeout} estricto que aborta
compilaciones colgadas o maliciosas.
\end{AlertaSeguridad}

\subsection{Procesamiento de \textit{insights} de reclutamiento mediante IA}
\label{subsec:seq_ia}

Los \textit{insights} del reclutador se generan invocando a \textbf{Gemini 2.5 Flash} a
través del servicio \comando{RecruiterAiService}. El modelo opera anclado al JSON de
perfiles que el backend inyecta en el \textit{prompt}, con configuración fija \\
(\texttt{maxOutputTokens=8192}, \texttt{temperature=0.7}, \texttt{thinkingBudget=0}) y
\textbf{cuatro modos} de operación seleccionados por el parámetro \texttt{mode}, descritos
en la tabla~\ref{tab:ia-modos}.

\vspace{0.2cm}
\noindent
\renewcommand{\arraystretch}{1.3}
\fontsize{8.5}{10.2}\selectfont\sffamily
\begin{tabularx}{\textwidth}{|L{3.4cm}|Y|}
\hline
\rowcolor{headerTabla}
\sffamily\textbf{Modo (\texttt{mode})} & \sffamily\textbf{Comportamiento del modelo} \\
\hline
\texttt{weekly\_insights} & Resume en un párrafo ($\le$60 palabras) el foco semanal del reclutador y sugiere una acción concreta. \\ \hline
\rowcolor{grisTecnico}
\texttt{compare\_verdict} & Evalúa cada perfil frente al rol: bloque Markdown con \textit{Pros}, \textit{Contras}, \textit{Riesgos} y una recomendación final. \\ \hline
\texttt{ranking\_justification} & Explica en exactamente 2 frases por qué un perfil encabeza el \textit{ranking}, citando \textit{skills} y proyectos reales. \\ \hline
\rowcolor{grisTecnico}
\texttt{nlp\_search} & Traduce una consulta en lenguaje natural a un objeto JSON de filtros estructurados (seniority, skills, institución, empresa). \\ \hline
\end{tabularx}
\label{tab:ia-modos}

\paragraph{Traducción de errores \textit{upstream}.} El punto crítico de diseño es que un
fallo de Gemini \textbf{no se propaga como un 500 genérico}. La factoría
\comando{AiServiceException.fromUpstream()} mapea el estado del proveedor a un código
accionable y un mensaje en español para la interfaz del reclutador, según la
tabla~\ref{tab:ia-errores}.

\vspace{0.2cm}
\noindent
\renewcommand{\arraystretch}{1.3}
\fontsize{9}{10.8}\selectfont\sffamily
\begin{tabularx}{\textwidth}{|C{2.6cm}|C{2.2cm}|Y|}
\hline
\rowcolor{headerTabla}
\sffamily\textbf{Estado Gemini} & \sffamily\textbf{Estado devuelto} & \sffamily\textbf{Mensaje al usuario} \\
\hline
429 & 429 & ``El asistente IA alcanzó su límite de uso. Intenta de nuevo más tarde.'' \\
\hline
\rowcolor{grisTecnico}
500 / 502 / 503 / 504 & 503 & ``El asistente IA no está disponible en este momento.'' \\
\hline
otros & 502 & ``Error del asistente IA (\textit{código}): \textit{mensaje}.'' \\
\hline
\end{tabularx}
\label{tab:ia-errores}

\begin{NotaTecnica}
Tanto el CV Studio como los \textit{insights} de reclutamiento reutilizan la misma
integración cruda con Gemini vía \texttt{HttpURLConnection} (\texttt{connectTimeout}
15\,s, \texttt{readTimeout} 120\,s). Lo que cambia entre ambos es el \textit{system
instruction} y el modo, no el transporte: una sola disciplina de manejo de errores
\textit{upstream} cubre toda la superficie de IA del sistema.
\end{NotaTecnica}

\subsection{Anclaje del asistente y modos del CV Studio}
\label{subsec:ia_cv_modos}

El asistente del CV Studio (\comando{GeminiAiService}) mantiene un \textbf{historial
multi-turno} de la conversación, de modo que el reescribir un párrafo conserva el contexto
de los intercambios previos. Su \textit{system instruction} \textbf{ancla} la salida a un
formato LaTeX seguro: según el modo, instruye al modelo a usar comandos LaTeX en modo
matemático para símbolos (\texttt{\$\textbackslash geq\$}, \texttt{\$\textbackslash star\$})
o a limitarse a texto ASCII y caracteres latinos acentuados, evitando viñetas Unicode que
romperían la compilación posterior. Esta restricción de formato es lo que permite que la
salida del modelo se compile directamente con Tectonic sin pasos de saneamiento adicionales
más allá de \comando{UnicodeLatexSanitizer}.

\begin{AvisoNota}
La diferencia de gobernanza entre ambos asistentes de IA es ilustrativa: el del reclutador
ancla la salida al \textbf{JSON de perfiles} para impedir alucinaciones sobre candidatos
inexistentes, mientras que el del CV Studio ancla la salida al \textbf{formato LaTeX} para
garantizar compilabilidad. En ambos casos, el \textit{system instruction} es el mecanismo
de control que mantiene al modelo dentro de límites operativos seguros.
\end{AvisoNota}
      % 4.3 Diagramas de secuencia
% =============================================================================
% 4.4 ARQUITECTURA DE COMUNICACIÓN EN TIEMPO REAL
% =============================================================================
\section{Arquitectura de Comunicación en Tiempo Real}
\label{sec:realtime}

La plataforma ofrece chat y notificaciones \textbf{en tiempo real} sin implementar un
servidor de \textit{WebSockets} dentro de Spring. Esta decisión —documentada como
ADR (sección~\ref{sec:adr}, pág.~\pageref{sec:adr})— delega la entrega \textit{push} en
\textbf{Supabase Realtime}, conservando la autoridad de validación en el backend.

\subsection{Infraestructura de mensajería en tiempo real}
\label{subsec:rt_infra}

El patrón es de \textbf{escritura validada por el backend y entrega \textit{push} por el
canal Realtime}. Combina lo mejor de dos mundos: la autoridad del servidor (validación,
idempotencia, persistencia) y la latencia mínima de la entrega directa al cliente
suscrito. La figura~\ref{fig:rt-arq} ilustra la topología.

\vspace{0.3cm}
\begin{figure}[h!]
\centering
\begin{tikzpicture}[node distance=1.0cm and 1.8cm,
    box/.style={rectangle, rounded corners=3pt, draw, font=\sffamily\scriptsize, align=center,
        minimum width=2.7cm, minimum height=0.95cm, inner sep=3pt},
    rel/.style={-{Stealth[length=2mm]}, semithick}]
    \node[box, fill=c4Persona!25] (emisor) {Emisor\\\scriptsize (SPA)};
    \node[box, fill=c4Componente!45, right=of emisor] (be) {ChatController\\\scriptsize valida + persiste};
    \node[box, fill=c4Datos!22, draw=c4Datos!70!black, right=of be] (db) {PostgreSQL\\\scriptsize core.chat\_messages};
    \node[box, fill=colorAcento!28, draw=colorAcento!80!black, below=of db] (rt) {Supabase\\Realtime};
    \node[box, fill=c4Persona!25, below=of emisor] (dest) {Destinatario\\\scriptsize (SPA suscrita)};

    \draw[rel] (emisor) -- node[c4etiqueta, above]{\scriptsize POST /messages} (be);
    \draw[rel] (be) -- node[c4etiqueta, above]{\scriptsize INSERT} (db);
    \draw[rel] (db) -- node[c4etiqueta, right]{\scriptsize WAL / trigger} (rt);
    \draw[rel] (rt) -- node[c4etiqueta, above]{\scriptsize push (sin polling)} (dest);
\end{tikzpicture}
\caption{Topología de mensajería: escritura validada por el backend, entrega \textit{push} por Supabase Realtime.}
\label{fig:rt-arq}
\end{figure}

El flujo concreto de un mensaje es:

\begin{enumerate}[noitemsep]
    \item El emisor envía el mensaje por REST a \comando{POST /api/v1/messages}.
    \item \comando{ChatController} valida que el llamante sea \textbf{participante} de la
    conversación y persiste el mensaje con \textbf{idempotencia} (clave de cliente,
    migración V7), evitando duplicados ante reintentos de red.
    \item \textbf{Supabase Realtime} detecta el cambio en la tabla mediante el registro
    de escritura anticipada (WAL) y lo entrega al destinatario suscrito, \textbf{sin
    \textit{polling}}.
\end{enumerate}

\begin{NotaTecnica}
La idempotencia es esencial en redes móviles inestables: si el cliente reintenta un
\texttt{POST} cuya respuesta se perdió, el backend reconoce la clave de idempotencia y
devuelve el mensaje ya persistido en lugar de crear un duplicado. Esta garantía se
implementa con una restricción \texttt{UNIQUE} parcial en la base de datos
(véase el capítulo~\ref{ch:basedatos}, pág.~\pageref{ch:basedatos}).
\end{NotaTecnica}

\subsection{Máquina de estados de una sesión de chat}
\label{subsec:rt_estados}

El ciclo de vida de un mensaje puede modelarse como una máquina de estados finitos. La
figura~\ref{fig:fsm-chat} la representa, incluyendo el \textbf{borrado lógico} que
preserva la integridad histórica (los mensajes no se eliminan físicamente, solo se marcan
con \texttt{deleted\_at}).

\vspace{0.3cm}
\begin{figure}[h!]
\centering
\begin{tikzpicture}[node distance=1.4cm and 2.0cm,
    estado/.style={rectangle, rounded corners=8pt, draw=c4Sistema!70!black, fill=c4Componente!30,
        font=\sffamily\scriptsize\bfseries, align=center, minimum width=2.2cm, minimum height=0.8cm},
    ini/.style={circle, fill=flujoInicio, minimum size=0.4cm, inner sep=0pt},
    rel/.style={-{Stealth[length=2mm]}, semithick}]
    \node[ini] (start) {};
    \node[estado, right=1.2cm of start] (compose) {Redactando};
    \node[estado, right=of compose] (sent) {Enviado\\(persistido)};
    \node[estado, right=of sent] (delivered) {Entregado\\(Realtime)};
    \node[estado, below=of delivered] (read) {Leído};
    \node[estado, below=of sent] (deleted) {Borrado lógico};

    \draw[rel] (start) -- (compose);
    \draw[rel] (compose) -- node[c4etiqueta, above]{POST /messages} (sent);
    \draw[rel] (sent) -- node[c4etiqueta, above]{push} (delivered);
    \draw[rel] (delivered) -- node[c4etiqueta, right]{visto} (read);
    \draw[rel] (sent) -- node[c4etiqueta, left]{eliminar} (deleted);
    \draw[rel] (delivered) -- node[c4etiqueta, above, sloped]{eliminar} (deleted);
\end{tikzpicture}
\caption{Máquina de estados de un mensaje de chat, con borrado lógico.}
\label{fig:fsm-chat}
\end{figure}

Las \textbf{notificaciones} siguen un patrón análogo y asíncrono: se crean en respuesta a
eventos (mensaje nuevo, \textit{like} recibido, etc.) y el usuario las consume vía
\comando{GET /api/v1/notifications}, marcándolas como leídas individualmente
(\comando{PATCH /\{id\}/read}) o en bloque (\comando{POST /read-all}). La tabla
\texttt{notifications} dispara, a su vez, eventos Realtime que actualizan el contador del
frontend sin recargar la página.

\subsection{Notificaciones por correo asíncronas (SMTP)}
\label{subsec:email}

Más allá de las notificaciones in-app, el backend envía \textbf{correos transaccionales}
mediante \comando{EmailService}, que usa el cliente de Spring Mail sobre SMTP
(\texttt{smtp.gmail.com:587} con STARTTLS). Los correos cubren los eventos de la
tabla~\ref{tab:emails}. El envío se realiza de forma que no bloquee la respuesta HTTP al
usuario, preservando la latencia percibida.

\vspace{0.2cm}
\noindent
\renewcommand{\arraystretch}{1.3}
\fontsize{8.5}{10.2}\selectfont\sffamily
\begin{tabularx}{\textwidth}{|L{3.6cm}|Y|}
\hline
\rowcolor{headerTabla}
\sffamily\textbf{Evento} & \sffamily\textbf{Correo enviado} \\
\hline
Recuperación de contraseña & Código OTP de un solo uso (sección~\ref{subsec:seq_otp}, pág.~\pageref{subsec:seq_otp}). \\ \hline
\rowcolor{grisTecnico}
Cambio de correo & Confirmación con enlace a \pathbreak{/api/v1/auth/confirm-email-change}. \\ \hline
Eventos administrativos & Comunicaciones del panel de administración (\pathbreak{/api/admin/email}). \\ \hline
\end{tabularx}
\label{tab:emails}

\subsection{Preocupaciones transversales del servidor}
\label{subsec:transversales}

Tres preocupaciones atraviesan todos los dominios del backend y se resuelven de forma
centralizada, no por dominio:

\begin{itemize}[noitemsep]
    \item \textbf{Idempotencia}: las operaciones de escritura sensibles a reintentos (envío
    de mensajes) aceptan una clave de idempotencia que evita duplicados, respaldada por una
    restricción de unicidad en la base de datos.
    \item \textbf{Registro estructurado} (\textit{logging}): cada servicio emite trazas con
    SLF4J; los errores no recuperables se registran con \textit{stack trace} completo en el
    servidor, mientras que al cliente solo viaja un mensaje seguro.
    \item \textbf{Zona horaria y serialización}: Jackson se configura globalmente en UTC y
    con \texttt{default-property-inclusion: NON\_NULL}, garantizando respuestas compactas y
    marcas de tiempo coherentes en toda la API.
\end{itemize}
        % 4.4 Comunicación en tiempo real
% =============================================================================
% 4.5 SEGURIDAD DEL SERVIDOR Y MANEJO DE EXCEPCIONES
% =============================================================================
\section{Seguridad del Servidor y Manejo de Excepciones}
\label{sec:excepciones}

La seguridad del backend se articula en dos planos complementarios: la \textbf{cadena de
filtros HTTP} que protege cada petición y la \textbf{traducción centralizada de
excepciones} que garantiza respuestas consistentes y seguras. Esta sección documenta
ambos, cerrando el descenso al nivel de código del componente servidor.

\subsection{Configuración de filtros HTTP}
\label{subsec:filtros}

La cadena de seguridad se define en \comando{SecurityConfig} como un
\textit{SecurityFilterChain} \textbf{stateless} \\
(\comando{SessionCreationPolicy.STATELESS}), con CSRF deshabilitado —no hay sesiones de
cookie tradicionales— y CORS gestionado explícitamente. Sobre la cadena estándar se
añaden filtros propios del paquete \texttt{config/}, resumidos en la
tabla~\ref{tab:filtros}.

\vspace{0.2cm}
\noindent
\renewcommand{\arraystretch}{1.3}
\fontsize{8.5}{10.2}\selectfont\sffamily
\begin{tabularx}{\textwidth}{|L{4.0cm}|Y|}
\hline
\rowcolor{headerTabla}
\sffamily\textbf{Filtro / Componente} & \sffamily\textbf{Función} \\
\hline
\texttt{CookieTokenFilter} & Permite leer el JWT desde una \textit{cookie} además del encabezado \texttt{Authorization}. \\
\hline
\rowcolor{grisTecnico}
\texttt{SecurityHeadersFilter} & Añade cabeceras de seguridad HTTP (endurecimiento de respuestas: \textit{nosniff}, \textit{frame-options}). \\
\hline
\texttt{NoCacheFilter} & Evita el cacheo de respuestas sensibles en navegador y proxies intermedios. \\
\hline
\rowcolor{grisTecnico}
\texttt{CorsConfig} & Define los orígenes permitidos (\texttt{CORS\_ALLOWED\_ORIGINS}) y los verbos/cabeceras admitidos. \\
\hline
\texttt{SpaFallbackFilter} & Reenvía las rutas de navegación de la SPA a \texttt{index.html} (servido por el mismo \textit{JAR}). \\
\hline
\end{tabularx}
\label{tab:filtros}

\paragraph{Rutas públicas.} La lista \comando{PUBLIC\_PATHS} declara los \textit{endpoints}
accesibles sin autenticación; \\ 
\comando{anyRequest().authenticated()} protege el resto. La
tabla~\ref{tab:public-paths} agrupa las rutas públicas reales por categoría.

\clearpage
\vspace{0.2cm}
\noindent
\renewcommand{\arraystretch}{1.3}
\fontsize{8}{9.6}\selectfont\sffamily
\begin{tabularx}{\textwidth}{|L{3.2cm}|Y|}
\hline
\rowcolor{headerTabla}
\sffamily\textbf{Categoría} & \sffamily\textbf{Rutas públicas} \\
\hline
Autenticación & \pathbreak{/api/auth/register}, \pathbreak{/api/auth/login}, \pathbreak{/api/auth/forgot-password/request}, \pathbreak{/api/auth/forgot-password/verify}, \pathbreak{/api/auth/forgot-password/reset}. \\
\hline
\rowcolor{grisTecnico}
API pública & \pathbreak{/api/ping}, \pathbreak{/api/v1/portfolio/public/**}, \pathbreak{/api/v1/projects/public/**}. \\
\hline
Infraestructura & \pathbreak{/error}, \pathbreak{/actuator/health}, \pathbreak{/swagger-ui/**}, \pathbreak{/swagger-ui.html}, \pathbreak{/v3/api-docs/**}. \\
\hline
\rowcolor{grisTecnico}
Estáticos SPA & \pathbreak{/}, \pathbreak{/index.html}, \pathbreak{/assets/**}, \pathbreak{/static/**}. \\
\hline
\end{tabularx}
\label{tab:public-paths}

\paragraph{Validación de firmas JWT.} El \comando{JwtDecoder} admite \textbf{dos modos}
según la configuración disponible, descritos en la tabla~\ref{tab:jwt-modos}. Al arranque,
\comando{warmupJwks()} pre-descarga las claves para evitar latencia en la primera
petición protegida.

\vspace{0.2cm}
\noindent
\renewcommand{\arraystretch}{1.3}
\fontsize{8.5}{10.2}\selectfont\sffamily
\begin{tabularx}{\textwidth}{|C{3.0cm}|C{2.6cm}|Y|}
\hline
\rowcolor{headerTabla}
\sffamily\textbf{Modo} & \sffamily\textbf{Algoritmo} & \sffamily\textbf{Fuente de la clave} \\
\hline
Secreto compartido & HS256 & \texttt{supabase.jwt-secret} (clave simétrica). \\ \hline
\rowcolor{grisTecnico}
JWKS remoto & RS256 / ES256 & \texttt{jwk-set-uri} (claves públicas de Supabase, cacheadas). \\ \hline
\end{tabularx}
\label{tab:jwt-modos}

\begin{AlertaSeguridad}
La derivación del rol se delega en \comando{AdminEmailPolicy.resolveRole()}, que
\textbf{degrada} cualquier intento de \texttt{role=admin} proveniente de un correo no
autorizado. Un atacante que manipule el \textit{claim} \texttt{role} de su token no puede
escalar a administrador: la autoridad efectiva se recalcula en el servidor a partir de la
lista de correos administrativos, nunca se confía ciegamente en el token.
\end{AlertaSeguridad}

\subsection{Jerarquía de excepciones técnicas y funcionales}
\label{subsec:jerarquia_exc}

La figura~\ref{fig:exc-jerarquia} muestra las excepciones propias y su mapeo a códigos
HTTP. La filosofía rectora es \textbf{no devolver nunca un 500 genérico} cuando el error
es recuperable o atribuible a una dependencia externa.

\vspace{0.3cm}
\begin{figure}[h!]
\centering
\begin{tikzpicture}[
    exc/.style={rectangle, draw=c4Sistema!70!black, fill=c4Componente!30, font=\sffamily\scriptsize,
        align=center, minimum width=3.4cm, minimum height=0.85cm, inner sep=3pt},
    http/.style={rectangle, rounded corners=8pt, fill=colorAcento!30, draw=colorAcento!80!black,
        font=\sffamily\scriptsize\bfseries, align=center, minimum width=2.2cm, minimum height=0.7cm},
    rel/.style={-{Stealth[length=2mm]}, semithick}]
    \node[exc, fill=c4Componente!55] (rt) at (0,2.6) {\texttt{RuntimeException}\\\scriptsize (raíz común)};
    \node[exc] (auth) at (-3.0,1.0) {\texttt{AuthException}\\\scriptsize (+ X-Blocked-Until)};
    \node[exc] (ai)   at (-3.0,-0.4) {\texttt{AiServiceException}\\\scriptsize fromUpstream()};
    \node[exc] (val)  at (-3.0,-1.8) {\texttt{MethodArgument}\\\texttt{NotValidException}};
    \node[http] (h401) at (3.0,1.0) {401 / 4xx};
    \node[http] (h5)   at (3.0,-0.4) {429 / 503 / 502};
    \node[http] (h400) at (3.0,-1.8) {400 + errores};

    \draw[rel] (auth.north) to[bend left=10] (rt.west);
    \draw[rel] (ai.east) to[bend right=20] (rt.south);
    \draw[rel, dashed] (auth) -- node[c4etiqueta, above]{mapea} (h401);
    \draw[rel, dashed] (ai) -- node[c4etiqueta, above]{mapea} (h5);
    \draw[rel, dashed] (val) -- node[c4etiqueta, above]{mapea} (h400);
\end{tikzpicture}
\caption{Jerarquía de excepciones del backend y su traducción a códigos HTTP.}
\label{fig:exc-jerarquia}
\end{figure}

\subsection{Estructura del \texttt{GlobalExceptionHandler} centralizado}
\label{subsec:global_handler}

El \comando{@RestControllerAdvice GlobalExceptionHandler} centraliza la traducción de
excepciones a envoltorios \comando{ApiResponse} (sección~\ref{subsec:apiresponse},
pág.~\pageref{subsec:apiresponse}), manteniendo los controladores limpios de lógica de
manejo de errores. La tabla~\ref{tab:handler} resume sus manejadores.

\vspace{0.2cm}
\noindent
\renewcommand{\arraystretch}{1.3}
\fontsize{8.5}{10.2}\selectfont\sffamily
\begin{tabularx}{\textwidth}{|L{4.6cm}|C{1.6cm}|Y|}
\hline
\rowcolor{headerTabla}
\sffamily\textbf{Excepción capturada} & \sffamily\textbf{HTTP} & \sffamily\textbf{Tratamiento} \\
\hline
\texttt{MethodArgumentNotValid} & 400 & Lista de errores por campo (\texttt{campo: mensaje}). \\
\hline
\rowcolor{grisTecnico}
\texttt{AuthException} & variable & Estado propio + cabecera \texttt{X-Blocked-Until} si aplica. \\
\hline
\texttt{IllegalArgumentException} & 400 & \texttt{badRequest} con el mensaje. \\
\hline
\rowcolor{grisTecnico}
\texttt{AuthorizationDenied} & 403 & ``Forbidden: insufficient role'' (no cae al 500). \\
\hline
\texttt{AiServiceException} & 429/503/502 & Estado propio del fallo \textit{upstream} de IA. \\
\hline
\rowcolor{grisTecnico}
\texttt{RuntimeException} / \texttt{Exception} & 500 & Último recurso, registrado en \textit{log} con traza completa. \\
\hline
\end{tabularx}
\label{tab:handler}

\begin{AlertaSeguridad}
El \textbf{orden} de los manejadores importa: \texttt{AuthorizationDeniedException} y
\texttt{AiServiceException} se capturan \textbf{antes} del manejador genérico de
\texttt{RuntimeException}, garantizando que las denegaciones de autorización (403) y los
fallos de IA (429/503/502) no se enmascaren como errores 500. Un 500 espurio no solo
degrada la experiencia, sino que puede filtrar detalles internos del \textit{stack}; por
ello el manejador genérico es el último recurso y registra la traza solo en el servidor.
\end{AlertaSeguridad}

\subsection{Síntesis de las capas de defensa}
\label{subsec:defensa_capas}

La seguridad del backend no descansa en un único control, sino en \textbf{capas de defensa
en profundidad} que se refuerzan mutuamente. La tabla~\ref{tab:defensa} las sintetiza, de
la más externa a la más interna.

\vspace{0.2cm}
\noindent
\renewcommand{\arraystretch}{1.3}
\fontsize{8.5}{10.2}\selectfont\sffamily
\begin{tabularx}{\textwidth}{|L{3.2cm}|Y|}
\hline
\rowcolor{headerTabla}
\sffamily\textbf{Capa} & \sffamily\textbf{Control} \\
\hline
Transporte & CORS, cabeceras de seguridad, \textit{no-cache} en respuestas sensibles. \\ \hline
\rowcolor{grisTecnico}
Autenticación & Verificación de JWT contra JWKS; sesión \textit{stateless}. \\ \hline
Autorización & Roles derivados en servidor; degradación de \texttt{admin} no autorizado. \\ \hline
\rowcolor{grisTecnico}
Validación & Bean Validation declarativa en los DTOs de entrada. \\ \hline
Datos & Políticas RLS en el motor (capítulo~\ref{ch:basedatos}, pág.~\pageref{ch:basedatos}). \\ \hline
\end{tabularx}
\label{tab:defensa}
       % 4.5 Seguridad y manejo de excepciones
% =============================================================================
% 4.6 CONFIGURACIÓN, EMPAQUETADO Y DESPLIEGUE DEL BACKEND
% =============================================================================
\section{Configuración, Empaquetado y Despliegue}
\label{sec:config_backend}

El diseño del backend no termina en el código: incluye su \textbf{configuración
externalizada}, su \textbf{empaquetado monolítico} y su \textbf{imagen de ejecución}. Esta
sección documenta esos tres aspectos, que la norma ISO/IEC/IEEE~15289:2024 clasifica como
parte del ítem de información de \textit{configuración del software}.

\subsection{Configuración por perfiles y variables de entorno}
\label{subsec:config_perfiles}

La configuración se externaliza con el mecanismo de \textbf{perfiles de Spring Boot}: un
fichero base \pathbreak{application.yml} y dos sobrecargas por entorno
(\pathbreak{application-local.yml} y \pathbreak{application-prod.yml}). El perfil activo se
selecciona con la variable \texttt{SPRING\_PROFILES\_ACTIVE} (por defecto \texttt{local}).
\textbf{Ningún secreto se escribe en el código}: todos llegan como variables de entorno,
como muestra la tabla~\ref{tab:config-vars}.

\vspace{0.2cm}
\noindent
\renewcommand{\arraystretch}{1.3}
\fontsize{8.5}{10.2}\selectfont\sffamily
\begin{tabularx}{\textwidth}{|L{4.2cm}|Y|}
\hline
\rowcolor{headerTabla}
\sffamily\textbf{Variable de entorno} & \sffamily\textbf{Propósito} \\
\hline
\texttt{SUPABASE\_URL} & URL base del proyecto Supabase (Auth, Storage, REST). \\ \hline
\rowcolor{grisTecnico}
\texttt{SUPABASE\_ANON\_KEY} & Clave pública para operaciones con RLS. \\ \hline
\texttt{SUPABASE\_SERVICE\_ROLE\_KEY} & Clave privilegiada del backend (evade RLS); nunca se expone. \\ \hline
\rowcolor{grisTecnico}
\texttt{GEMINI\_API\_KEY} & Clave del motor de IA Gemini 2.5 Flash. \\ \hline
\texttt{CORS\_ALLOWED\_ORIGINS} & Orígenes autorizados para peticiones \textit{cross-origin}. \\ \hline
\rowcolor{grisTecnico}
\texttt{APP\_BASE\_URL} · \texttt{APP\_FRONTEND\_URL} & URLs absolutas para enlaces de correo y \textit{callbacks} OAuth. \\ \hline
\end{tabularx}
\label{tab:config-vars}

\begin{lstlisting}[language=yaml, caption={Extracto de configuración externalizada (\texttt{application.yml}).}, label={lst:appyml}]
spring:
  profiles:
    active: ${SPRING_PROFILES_ACTIVE:local}
  servlet:
    multipart:
      max-file-size: 20MB        # limite de subida de archivos
      max-request-size: 25MB
supabase:
  url: ${SUPABASE_URL:}
  service-role-key: ${SUPABASE_SERVICE_ROLE_KEY:}
gemini:
  api-key: ${GEMINI_API_KEY:}
\end{lstlisting}

\begin{AlertaSeguridad}
La separación entre \texttt{SUPABASE\_ANON\_KEY} (pública) y
\texttt{SUPABASE\_SERVICE\_ROLE\_KEY} (privilegiada) es la piedra angular del modelo de
seguridad de datos. La clave de servicio evade las políticas RLS del motor
(capítulo~\ref{ch:basedatos}, sección~\ref{sec:rls}, pág.~\pageref{sec:rls}) y, por tanto,
vive \textbf{únicamente} como variable de entorno del backend, jamás en el \textit{bundle}
del cliente.
\end{AlertaSeguridad}

\subsection{Empaquetado monolítico: un único \textit{JAR}}
\label{subsec:empaquetado}

EthosHub se distribuye como un \textbf{monolito desplegable}: un solo artefacto \textit{JAR}
que contiene tanto el backend Spring como el \textit{bundle} estático del frontend. El
\textit{script} \comando{build.sh} orquesta el proceso, ilustrado en la
figura~\ref{fig:build-monolito}.

\vspace{0.3cm}
\begin{figure}[h!]
\centering
\begin{tikzpicture}[node distance=0.6cm and 1.2cm,
    paso/.style={rectangle, rounded corners=3pt, draw=c4Sistema!70!black, fill=c4Componente!30,
        font=\sffamily\scriptsize, align=center, minimum width=2.7cm, minimum height=0.85cm, inner sep=3pt},
    art/.style={paso, fill=c4Datos!22, draw=c4Datos!70!black},
    rel/.style={-{Stealth[length=2mm]}, semithick}]
    \node[paso] (fe) {npm run build\\\scriptsize Vite + tsc};
    \node[art, right=of fe] (dist) {\pathbreak{dist/}\\\scriptsize bundle SPA};
    \node[paso, right=of dist] (cp) {copia a\\\pathbreak{static/}};
    \node[paso, below=of cp] (mvn) {mvn package\\\scriptsize Spring Boot};
    \node[art, left=of mvn] (jar) {\texttt{app.jar}\\\scriptsize monolito};

    \draw[rel] (fe) -- (dist);
    \draw[rel] (dist) -- (cp);
    \draw[rel] (cp) -- (mvn);
    \draw[rel] (mvn) -- (jar);
\end{tikzpicture}
\caption{Empaquetado monolítico: el \textit{bundle} de la SPA se incrusta en el \textit{JAR} del backend.}
\label{fig:build-monolito}
\end{figure}

El proceso consta de cuatro pasos verificados por el \textit{script}:

\begin{enumerate}[noitemsep]
    \item Compilación del frontend (\comando{npm run build}), que ejecuta \texttt{tsc}
    (verificación de tipos) y \texttt{vite build} (empaquetado optimizado a \pathbreak{dist/}).
    \item Verificación de que \pathbreak{dist/} existe (detecta fallos silenciosos del build).
    \item Copia de \pathbreak{dist/} al directorio \pathbreak{src/main/resources/static/} del
    backend, previa limpieza del contenido anterior.
    \item Empaquetado con Maven (\comando{mvn package}), que produce un \textit{JAR}
    ejecutable que sirve la SPA y la API bajo el mismo origen.
\end{enumerate}

\begin{NotaTecnica}
El empaquetado de origen único (\textit{same-origin}) elimina por completo la clase de
problemas CORS en producción y simplifica el despliegue a un solo proceso. El \textit{flag}
\texttt{-{}-skip-fe} permite reconstruir solo el backend durante el desarrollo, acelerando
las iteraciones que no tocan el cliente.
\end{NotaTecnica}

\subsection{Imagen de ejecución con compilador LaTeX (Tectonic)}
\label{subsec:docker}

La imagen de contenedor se construye en \textbf{dos etapas} (\textit{multi-stage build}):
una etapa de compilación con el JDK~17 completo y una etapa de ejecución con solo el JRE.
La etapa de ejecución instala \textbf{Tectonic 0.15.0} (binario estático \textit{musl},
sin dependencias adicionales) y \textbf{precalienta su caché de paquetes} compilando un
documento de prueba (\texttt{warmup.tex}); así, la primera compilación de CV de un usuario
real no paga la descarga de paquetes LaTeX. La tabla~\ref{tab:docker} resume las etapas.

\vspace{0.2cm}
\noindent
\renewcommand{\arraystretch}{1.3}
\fontsize{8.5}{10.2}\selectfont\sffamily
\begin{tabularx}{\textwidth}{|L{3.4cm}|Y|}
\hline
\rowcolor{headerTabla}
\sffamily\textbf{Etapa} & \sffamily\textbf{Contenido} \\
\hline
\texttt{build} (JDK 17) & Compila el proyecto Maven y produce \texttt{app.jar}. \\ \hline
\rowcolor{grisTecnico}
\texttt{runtime} (JRE 17) & Instala Tectonic, precalienta su caché y ejecuta \texttt{java -jar app.jar} en el puerto 8080. \\ \hline
\end{tabularx}
\label{tab:docker}

Esta imagen es la que materializa la compilación de currículums descrita en la
secuencia de la figura~\ref{fig:seq-cv} (pág.~\pageref{fig:seq-cv}): el binario LaTeX
preinstalado es el ``proceso externo'' que transforma el código del CV en un PDF
descargable.

\subsection{Configuración del cliente HTTP saliente (\texttt{RestClient})}
\label{subsec:restclient}

Las llamadas \textbf{salientes} del backend hacia servicios externos (Supabase Storage,
entre otros) se canalizan por instancias de \comando{RestClient} de Spring, configuradas en
\comando{RestClientConfig}. El \comando{SupabaseStorageService}, por ejemplo, construye su
cliente con la URL base de Storage y las cabeceras de autorización
(\texttt{Authorization: Bearer <service\_role>} y \texttt{apikey}) preestablecidas, de modo
que cada operación de subida o borrado ya viaja autenticada. La tabla~\ref{tab:salientes}
resume las integraciones salientes y su transporte.

\vspace{0.2cm}
\noindent
\renewcommand{\arraystretch}{1.3}
\fontsize{8.5}{10.2}\selectfont\sffamily
\begin{tabularx}{\textwidth}{|L{3.4cm}|L{3.0cm}|Y|}
\hline
\rowcolor{headerTabla}
\sffamily\textbf{Destino} & \sffamily\textbf{Transporte} & \sffamily\textbf{Uso} \\
\hline
Supabase Storage & \texttt{RestClient} & Subida/borrado en \textit{bucket} \texttt{profile-assets}. \\ \hline
\rowcolor{grisTecnico}
Gemini API & \texttt{HttpURLConnection} & Asistencia IA del CV e \textit{insights} de reclutamiento. \\ \hline
SMTP (Gmail) & Spring Mail & Correos transaccionales (OTP, confirmaciones). \\ \hline
\rowcolor{grisTecnico}
Supabase JWKS & \texttt{JwtDecoder} & Claves públicas para validar JWT. \\ \hline
\end{tabularx}
\label{tab:salientes}

\bigskip
Con esto se cierra el diseño detallado del backend. El capítulo~\ref{ch:frontend}
(pág.~\pageref{ch:frontend}) documenta el cliente que consume esta API, incluyendo cómo el
interceptor de Axios reacciona precisamente a los códigos 401/403 aquí definidos, y cómo
la SPA se incrusta en el mismo \textit{JAR} monolítico que aquí se ha descrito.
   % 4.6 Configuración, empaquetado y despliegue
